\sect{Приложение А. Алгоритмы Forward и Viterbi}{appendix-hmm}
%
В настоящем приложении приведены полные рекуррентные формулы и
псевдокод классических алгоритмов Forward и Viterbi для скрытой
марковской модели, используемой в~\Secref{background}. Для
самостоятельности изложения здесь повторяется обозначение
$\lambda=(A,B,\pi)$, где $A=(A_{ij})$~--- матрица переходов согласно
\eqref{eq:transition-matrix}, $b_i(o)$~--- эмиссионная плотность,
$\pi=(\pi_i)$~--- начальное распределение.

Алгоритм Forward вычисляет полную вероятность наблюдаемой
последовательности $O=(o_1,\dots,o_T)$ при модели $\lambda$ через
вспомогательные величины
$\alpha_t(i) = \Prb(o_1,\dots,o_t,\,s_t=i\mid\lambda)$. Рекуррентная
формула имеет вид
\begin{equation}
\alpha_t(j) = b_j(o_t)\sum_{i=1}^{N}\alpha_{t-1}(i)\,A_{ij},
\eqlab{forward-app}
\end{equation}
с инициализацией $\alpha_1(j)=\pi_j\,b_j(o_1)$. Полная вероятность
наблюдений равна $\Prb(O\mid\lambda)=\sum_{j=1}^{N}\alpha_T(j)$. На
практике для длинных последовательностей применяется покадровая
нормировка $\tilde\alpha_t(j)=\alpha_t(j)\bigl/\sum_k\alpha_t(k)$,
что устраняет численное переполнение и одновременно даёт
апостериорную оценку $\Prb(s_t=j\mid o_1,\dots,o_t)$, используемую
для онлайн-выдачи $\hat\varphi(t)=\argmaxop_j\tilde\alpha_t(j)$.

\par\addvspace{2pt}
{\centering
\begin{minipage}{0.97\linewidth}
\footnotesize
\textbf{Алгоритм 1.} Forward-рекурсия для онлайн-трекера. \\[1pt]
\rule{\linewidth}{0.4pt}\vspace{1pt}\\
\textbf{Вход:} партитура $S=(p_1,\dots,p_N)$, матрица $A$,
эмиссионный параметр $\sigma$. \\
1: $\alpha\gets(1,0,\dots,0)\in\Real^N$ \\
2: \textbf{для} каждой поступающей ноты $o$ \textbf{выполнить} \\
3: \quad \textbf{для} $j\gets 1,\dots,N$ \textbf{выполнить} \\
4: \quad\quad $b_j\gets\mathcal{N}(o;\,p_j,\,\sigma^2)$ \\
5: \quad\quad $\alpha'_j\gets b_j\sum_{i=1}^{N}\alpha_i\,A_{ij}$ \\
6: \quad \textbf{конец для} \\
7: \quad $\alpha\gets\alpha'\bigl/\sum_k\alpha'_k$ \\
8: \quad \textbf{выдать} $\hat\varphi\gets\argmaxop_j\alpha_j$ \\
9: \textbf{конец для} \\
\rule{\linewidth}{0.4pt}
\end{minipage}\par}
\addvspace{4pt}

Алгоритм Viterbi возвращает не апостериорное распределение, а одну
наиболее вероятную траекторию состояний $s_1^*,\dots,s_T^*$. Он
оперирует величинами
$\delta_t(i) = \max_{s_1,\dots,s_{t-1}}
\Prb(s_1,\dots,s_t=i,\,o_1,\dots,o_t\mid\lambda)$
и обратными указателями $\psi_t(i)$:
\begin{equation}
\eqlab{viterbi-delta}
\delta_t(j) = b_j(o_t)\,\max_{i}\bigl[\delta_{t-1}(i)\,A_{ij}\bigr],
\end{equation}
\begin{equation}
\eqlab{viterbi-psi}
\psi_t(j) = \argmaxop_{i}\bigl[\delta_{t-1}(i)\,A_{ij}\bigr],
\end{equation}
с инициализацией $\delta_1(j)=\pi_j\,b_j(o_1)$ и $\psi_1(j)=0$.
Оптимальная траектория восстанавливается обратным проходом:
$s_T^*=\argmaxop_j\delta_T(j)$, $s_{t-1}^*=\psi_t(s_t^*)$.

В нашем трекере Viterbi используется в \textit{режиме отложенного
анализа}: после исполнения каждого фрагмента партитуры (например,
после двойной чёрточки) выполняется обратный проход по
накопленным $\delta$ и $\psi$, что позволяет уточнить выровнянную
последовательность для дальнейшего обучения CNN-эмиссии и
обновления статистик переходов на следующих исполнениях.

Сложность одного шага обоих алгоритмов составляет $O(N)$ за счёт
полосной структуры матрицы $A$ из~\eqref{eq:transition-matrix}:
вместо суммирования по всем состояниям достаточно обработать три
ненулевых перехода. Память~--- $O(N)$ для Forward и $O(N T)$ для
Viterbi (требуется хранить указатели $\psi_t$ за весь интервал
анализа).
